\section{Guide d'utilisation du programme}

Pour exécuter le programme, ouvrez un terminal dans le répertoire du projet et lancez :

\begin{verbatim}
python -m src.main
\end{verbatim}

Le programme dispose de deux modes principaux : un mode commande classique et un mode interactif.

\subsection{Mode commande classique}

Pour obtenir l'aide générale et les options disponibles, utilisez :

\begin{verbatim}
python -m src.main -h
\end{verbatim}

La sortie affiche :

\begin{verbatim}
usage: Graphoni [-h] [-a {kruskal,prim,reverse}] [-s] [-i] [-v] [input] [output]

Util for finding MST of graph

positional arguments:
  input                 input file
  output                output file

options:
  -h, --help            show this help message and exit
  -a, --algorithm {kruskal,prim,reverse}
                        choose algorithm
  -s, --second          find second MST
  -i, --interactive     run interactive mode
  -v, --version         show version
\end{verbatim}

\textbf{Explications des options :}
\begin{itemize}
    \item \texttt{-a, --algorithm} : choisir l'algorithme pour le calcul de l'arbre couvrant minimal (Kruskal, Prim ou Reverse Delete). Par défaut : Kruskal.
    \item \texttt{-s, --second} : calculer le second MST.
    \item \texttt{-i, --interactive} : lancer le mode interactif.
    \item \texttt{input} : fichier d'entrée contenant le graphe.
    \item \texttt{output} : fichier de sortie pour enregistrer le résultat.
\end{itemize}

\subsection{Mode interactif}

Le mode interactif permet de modifier le graphe et de calculer le MST en temps réel.
Pour l'activer, lancez :

\begin{verbatim}
python -m src.main -i
\end{verbatim}

Les principales commandes disponibles :

\begin{itemize}
    \item \texttt{load <path>} : charger un graphe depuis un fichier.
    \item \texttt{save <path>} : enregistrer le graphe ou le MST dans un fichier.
    \item \texttt{add <u> (<v> <weight>)} : ajouter un sommet ou une arête au graphe.
    \item \texttt{remove <u> (<v>)} : supprimer un sommet ou une arête.
    \item \texttt{mst (<algorithm> -s)} : calculer le MST. On peut spécifier l'algorithme et l'option step-by-step.
    \item \texttt{update <u> <v> (<weight>)} : mettre à jour le graphe et recalculer le MST après modification.
    \item \texttt{second} : trouver le second MST.
    \item \texttt{keep} : remplacer le graphe original par le graphe résultat pour les prochaines opérations.
    \item \texttt{origin} : afficher le graphe original.
    \item \texttt{repos} : repositionner le graphe.
    \item \texttt{help} : afficher l'aide dans le mode interactif.
    \item \texttt{exit} : quitter le mode interactif.
\end{itemize}

\subsection{Remarques pratiques}

\begin{itemize}
    \item Toujours lancer le programme depuis le répertoire principal du projet pour que les chemins relatifs fonctionnent correctement.
    \item Le mode interactif est particulièrement utile pour tester différentes modifications sur le graphe sans avoir à relancer le programme.
\end{itemize}
